\documentclass[8pt,a4paper]{altacv}

% ---------- Page geometry: maximize usable area on A4 ----------
\geometry{left=0.60cm,right=0.60cm,top=0.55cm,bottom=0.55cm,columnsep=1.05cm}
\usepackage[utf8]{inputenc}
\usepackage[T1]{fontenc}
\usepackage[default]{lato}
\usepackage{paracol}
\columnratio{0.31}
\usepackage{enumitem}
\setlist[itemize]{leftmargin=*,itemsep=2.2pt,topsep=1.2pt,parsep=0pt,partopsep=0pt}
\setlength{\parskip}{1pt}
\renewcommand{\itemmarker}{{\scriptsize\textbullet}}

% ---------- Refined medical palette ----------
\definecolor{MainBlue}{HTML}{087E8B}
\definecolor{DeepNavy}{HTML}{12324A}
\definecolor{SlateGrey}{HTML}{314652}
\definecolor{Ink}{HTML}{202B33}
\definecolor{Sidebar}{HTML}{F1F7F8}
\colorlet{heading}{DeepNavy}
\colorlet{accent}{MainBlue}
\colorlet{emphasis}{SlateGrey}
\colorlet{body}{Ink}

\name{Wafae Youssfi}
\tagline{Ingénieure biomédicale double diplômée dotée d’une bonne compréhension des environnements hospitaliers et des technologies médicales, avec un fort intérêt pour les solutions de santé innovantes et leur déploiement dans les environnements cliniques}
\personalinfo{}
\photo{2.8cm}{picture5}

\begin{document}
\makecvheader
\begin{paracol}{2}
\backgroundcolor{c[0]}{Sidebar}

% ================= LEFT COLUMN =================
\cvsection{CONTACT}
\small
\vspace{2pt}
\email{wafaey020@gmail.com}
\vspace{2pt}
\phone{+33 781753210}
\vspace{2pt}
\location{Rue Paul Collomp , Clermont-Ferrand, France}
\vspace{2pt}
\location{Rue Saint Denis , Sannois (Région Parisienne), France}
\normalsize
\medskip

\cvsection{COMPÉTENCES}
\textbf{\color{accent}{Hard skills}}
\begin{itemize}
  \item Microsoft Pack Office
  \item Python, C, MATLAB
  \item Visual Studio Code, Jupyter Notebook, ArgoUML, Paravision (MRI Software)
  \item Customer Relationship Management (CRM)
\end{itemize}
\medskip
\textbf{\color{accent}{Soft skills}}
\begin{itemize}
  \item Sens d’analyse
  \item Gestion du stress
  \item Travail en équipe
  \item Négociation
  \item Adaptabilité
  \item Gestion de projet
  \item Esprit d'entrepreneuriat
\end{itemize}
\medskip

\cvsection{CERTIFICATIONS}
\begin{itemize}
  \item Outils technologiques d’IA
  \item Gestion de Maintenance Assistée par Ordinateur (GMAO)
  \item Techniques de ventes commerciales
  \item Secourisme
  \item Anglais médical
\end{itemize}
\medskip

\cvsection{LANGUES}
\begin{itemize}
  \item Arabe : Maternelle
  \item Français : C2
  \item Anglais : C2
  \item Turc : B1--B2
\end{itemize}
\medskip

\cvsection{CENTRES D'INTÉRÊT}
\begin{itemize}
  \item Lecture
  \item Sport
  \item Chant
  \item Apprentissage de nouvelles compétences
  \item Ex-membre de l'Association Américaine Internationale des Femmes (AIWA)
\end{itemize}

% ================= RIGHT COLUMN =================
\switchcolumn
\cvsection{EXPÉRIENCES PROFESSIONNELLES}
\cvevent{\textbf{Stage de Recherche en Imagerie}}{INRAE -- Institut National de Recherche d'Agronomie et d'Environnement}{Avril 2025 -- Août 2025 (5 mois)}{Saint-Genès-Champanelle, France}
\begin{itemize}
  \setlength{\itemsep}{2.8pt}
  \setlength{\topsep}{1.2pt}
  \item Mission: Développement et validation de méthodes d’IRM métabolique en utilisant des agents hyperpolarisés transportables
  \item Réalisation d’une veille bibliographique approfondie sur les protocoles d’imagerie métabolique rapide en IRM hyperpolarisée, permettant l’acquisition de l’information métabolique
  \item Mise en oeuvre expérimentale de la séquence d’IRM \textit{IDEAL SPIRAL} sur des produits métaboliques
  \item Acquisition de données brutes d'IRM issues du logiciel de Paravision et leur post-traitement à l’aide de MATLAB pour la reconstruction des images métaboliques
  \item Étude de benchmark consistant à mesurer et analyser des indicateurs de qualité des images
  \item Exploration d’une autre approche d'IRM basée sur la génération d’impulsions spectrales-spatiales
  \item Participation active à des réunions stratégiques avec l’équipe de recherche, dédiées aux échanges scientifiques et au suivi de l’avancement des projets
\end{itemize}
\vspace{2pt}
\cvevent{\textbf{Stage Industriel -- Département Commercial \& Marketing}}{Mindray Medical Morocco}{Février 2023 -- Septembre 2023 (6 mois)}{Casablanca, Maroc}
\begin{itemize}
  \item Mission: Conception et développement d’une plateforme digitale de gestion dédiée aux respirateurs d’anesthésie de Mindray.
  \item Visite hospitalière pour appréhender le fonctionnement des machines d'anesthésie et leurs technicités
  \item Conduite d’une étude analytique technico-fonctionnelle de la gamme de produits de Mindray
  \item Participation à des congrès nationaux et internationaux dédiés au secteur de la santé
  \item Accompagnement et soutien de l’équipe technico-commerciale lors du cycle de vente des équipements
  \item Identification des besoins et des défis des professionnels de santé
\end{itemize}
\vspace{2pt}
\cvevent{\textbf{Stage Technique}}{CHU -- Hôpital des Spécialités}{Juillet 2022 -- Septembre 2022 (2 mois)}{Rabat, Maroc}
\begin{itemize}
  \item Réalisation d’un inventaire des équipements biomédicaux tout en développant une connaissance pratique du parc matériel hospitalier
  \item Coopération avec l’équipe technique dans le diagnostic des pannes d’équipements
  \item Contribution à la maintenance corrective des dispositifs du bloc opératoire de l’hôpital
  \item Assistance à des interventions chirurgicales (En orthopédie, neurochirurgie et ophtalmologie)
\end{itemize}
\vspace{2pt}
\cvevent{\textbf{Stage d'Apprentissage}}{Centre de Radiologie médicale}{Août 2021 (1 mois)}{Rabat, Maroc}
\begin{itemize}
  \item Assistance à la maintenance préventive d’équipement d’imagerie
  \item Assistance aux consultations médicales et formation au bon usage des appareils radiologiques
\end{itemize}
\vspace{2pt}
\cvevent{\textbf{Stage Hospitalier d’Initiation}}{Hôpital Militaire d’Instruction Mohammed V}{Juillet 2021 (1 mois)}{Rabat, Maroc}
\begin{itemize}
  \item Découverte et prise en main de différents dispositifs au sein de divers services hospitaliers
  \item Engagement dans le processus de maintenance corrective en collaboration avec l’équipe biomédicale
\end{itemize}

\cvsection{FORMATIONS}
\cvevent{Diplôme de Master en Ingénierie de Santé - Parcours Imagerie Médicale, Signaux et Systèmes}{École Centrale de Lyon -- Université de Claude Bernard (Polytech Lyon)}{Septembre 2024 -- Novembre 2025}{Lyon, France}
\cvevent{Diplôme d’Ingénieur d’État en Génie Biomédical}{Université Mohammed VI des Sciences et de la Santé - École Supérieure Mohammed VI d’Ingénieurs en Sciences de la Santé}{Septembre 2020 -- Octobre 2023}{Casablanca, Maroc}
\cvevent{Diplôme de DEUG en Sciences de la Matière Physique}{Université Mohammed V de Rabat -- Faculté des Sciences de Rabat}{Septembre 2017 -- Septembre 2020}{Rabat, Maroc}
\end{paracol}
\end{document}
